\chapter{Dollar-Cost Average the Right Thing}

The core method of successful investing is almost offensively simple:

\begin{quote}
Buy low, sell high.
\end{quote}

That is all. No sentence is more central. No technical system is more important. No complicated vocabulary can replace it.

Yet this is one of those secrets everybody knows and almost nobody can use. The secret is not in the sentence itself. The secret is in the doing.

Many important truths have this shape. Growth is also simple: improve a little every day. If every day is too much, improve a little every week. Even that is enough to create compounding over time. Everyone understands this. Very few people live this way.

Why?

Because reality usually refuses to provide immediate feedback. One unit of effort today may not return one unit of result today. It may return next week, next month, next year, or much later. The person who requires instant confirmation will abandon many correct actions before their effect appears.

Investing is the same. The sentence ``buy low, sell high'' is clear, but each word hides a difficult judgment.

\begin{enumerate}
\item ``Low'' and ``high'' are relative values, not absolute prices. A stock at \(100\) may be cheap if the company is worth far more. A stock at \(5\) may be expensive if the company is collapsing.
\item ``Buy'' and ``sell'' are not necessarily total actions. How much should be bought? How much should be sold? What proportion of total capital should be used?
\item The reference point is value, but value is hard to calculate. Different people use different theories, different formulas, and different assumptions.
\item In the end, the investor must choose a method and bear the consequence of being wrong.
\end{enumerate}

This is where many people become dependent. They want someone else to give the answer. They do not want to study, compare, calculate, or decide. They want a conclusion without paying the price of judgment.

That habit is expensive everywhere, but in investing it is fatal. In ordinary life, depending on others may merely annoy people. In the capital market, depending on others without understanding makes a person prey.

A qualified investor must gradually leave the habit of asking for ready-made answers. Other people's explanations, interpretations, and choices belong to them. Your capital belongs to you. Your outcome belongs to you. Therefore your homework must also belong to you.

\section*{The Simplest Useful Strategy}

There is, however, one strategy simple enough for nearly everyone to understand and useful enough for nearly everyone to study seriously:

\begin{quote}
Dollar-cost averaging: regularly invest the same amount of money into one or several growing assets.
\end{quote}

In this chapter we will use the phrase ``fixed-period, fixed-amount buying'' and the common English phrase ``dollar-cost averaging'' for the same idea.

The procedure is plain.

\begin{enumerate}
\item Choose an asset after study, preferably a company or basket of companies with strong long-term growth.
\item Set a fixed interval: weekly, monthly, quarterly, or another regular period.
\item At each interval, invest the same amount of money.
\item Ignore short-term price movement when executing the scheduled purchase.
\end{enumerate}

Suppose an investor puts \(1\) unit of capital into the same stock each month for three months.

\begin{center}
\begin{adjustbox}{max width=\linewidth}
\begin{tabular}{lccc}
\hline
Month & Amount invested & Stock price & Shares bought \\
\hline
January & \(1\) & \(10\) & \(0.100\) \\
February & \(1\) & \(15\) & \(0.067\) \\
March & \(1\) & \(13\) & \(0.077\) \\
\hline
Total / result & \(3\) & \(12.32\) & \(0.244\) \\
\hline
\end{tabular}
\end{adjustbox}
\end{center}

The simple average of the three monthly prices is:

\[
\frac{10+15+13}{3}=12.67
\]

But the investor's average purchase cost is about:

\[
\frac{3}{0.24359}=12.32
\]

The exact number depends on unrounded shares, but the lesson is clear. Because the investor spends the same amount each period, more shares are bought when the price is low and fewer shares are bought when the price is high. The method does not magically predict the low point. It captures an average across time while automatically giving heavier weight to lower prices.

That is why dollar-cost averaging is a risk-reduction tool. At the moment of purchase, the investor usually cannot know with certainty whether the current price is a short-term low, whether it will fall further next week, or whether it will rise immediately. Fixed-period, fixed-amount buying removes the need to answer that impossible short-term question every time.

\section*{What It Saves You From}

Dollar-cost averaging is powerful partly because of what it prevents.

It prevents the beginner from staring at price charts all day and pretending that anxiety is analysis. It prevents the habit of waiting forever for the perfect entry point. It prevents the repeated emotional mistake of buying only after excitement has already pushed the price up. It prevents the opposite mistake of freezing when a lower price appears, even though the lower price was supposedly what the investor wanted.

The strategy can be summarized as a small machine:

\begin{center}
\begin{adjustbox}{max width=\linewidth}
\begin{tikzpicture}[node distance=1.5cm, every node/.style={font=\small}]
\node[draw, rounded corners, align=center, minimum width=2.7cm, minimum height=0.85cm] (study) {Study first};
\node[draw, rounded corners, align=center, minimum width=2.7cm, minimum height=0.85cm, right=of study] (select) {Select growth};
\node[draw, rounded corners, align=center, minimum width=2.7cm, minimum height=0.85cm, right=of select] (schedule) {Set schedule};
\node[draw, rounded corners, align=center, minimum width=2.7cm, minimum height=0.85cm, below=of schedule] (buy) {Buy fixed amount};
\node[draw, rounded corners, align=center, minimum width=2.7cm, minimum height=0.85cm, left=of buy] (review) {Review thesis};
\draw[->] (study) -- (select);
\draw[->] (select) -- (schedule);
\draw[->] (schedule) -- (buy);
\draw[->] (buy) -- (review);
\draw[->] (review) -- (select);
\end{tikzpicture}
\end{adjustbox}
\end{center}

Notice the first box. It is not ``buy.'' It is ``study first.''

Many people misunderstand the strategy because they focus only on the easy part: the fixed date and fixed amount. Those are useful, but they are not the deepest source of return. The deepest source is the quality of the asset being bought.

A more honest formula is:

\[
\text{Investment result}
=
\text{growth of the asset}
\times
\text{effect of the strategy}.
\]

If the asset has no growth, the strategy has little to multiply. If the asset is declining for structural reasons, dollar-cost averaging may simply make the investor buy more of a worsening thing.

This is why the method is safest only when applied to something that has a genuine chance of long-term growth. The strategy reduces timing risk. It does not remove selection risk.

\section*{Selection Comes Before Execution}

The harder question is not whether one should buy regularly. The harder question is what deserves to be bought regularly.

This is the part nobody can outsource. The whole world of investors is trying to answer it. There is no single formula that can settle it forever. Even when a person has a theory, that theory must be tested, revised, and improved through time.

Still, the direction is clear. Before investing, the investor must study the attributes of growth.

What kind of company can keep increasing value? What creates its advantage? Is the market expanding? Is the product improving? Are customers becoming more dependent on it? Can the company reinvest capital at attractive rates? Does management allocate capital intelligently? Are the numbers consistent with the story? Are the risks visible, or merely hidden by enthusiasm?

These questions are not decorative. They are the homework.

Most people learned in school to do homework for someone else. They completed assignments to submit them, to satisfy a teacher, or to avoid punishment. Investing reverses the meaning of homework. Here the homework is done for yourself. The grade is not written in red ink. It is written in your capital account five or ten years later.

All investment homework should be done before the investment, not after. Before buying, one should have thought about what to buy, why to buy it, how much to buy, how to continue buying, under what conditions the thesis is broken, and under what conditions selling becomes reasonable.

Buying first and studying later is not investing. It is panic in slow motion.

\section*{Choice Is More Important Than Speed}

Earlier we distinguished doing the right thing from doing things efficiently. This distinction becomes sharper here.

If the thing itself is right, doing it slowly can still create positive accumulation. If the thing itself is wrong, greater speed only increases damage.

Dollar-cost averaging is a disciplined method, but discipline in the wrong direction is still wrong. A person who buys a poor asset every month is not prudent merely because the schedule is regular. He is systematically transferring capital into a bad choice.

This is why choice is more important than effort when one must rank them. In real life both are necessary. A person must choose and must act. But if forced to choose the more fundamental one, choose choice.

Conceptual clarity exists for this reason. We have spent many chapters polishing concepts because vague concepts produce poor choices. A person who thinks ``cheap'' means low nominal price may buy disaster. A person who thinks ``safe'' means emotionally comfortable may buy hidden risk. A person who thinks ``persistence'' means never revising a wrong idea may confuse stubbornness with strength.

Freedom comes only after the ability to choose correctly has been trained.

\section*{The Attention Version}

This strategy is not limited to money.

A person can dollar-cost average attention.

Imagine two readers who paid for the same course. One opens every lesson, reads carefully, thinks, writes notes, and acts on the exercises. The other pays, feels briefly satisfied, and rarely reads. Their purchase price is the same, but their cost per useful idea is completely different.

The first reader invests fixed attention at fixed intervals. The second reader only made a payment.

This is one of the book's central themes: attention is the most important asset under personal control. Money invested regularly can compound only after attention has selected the right target. Attention invested regularly can create the judgment that later protects and multiplies money.

The same rhythm applies to learning financial statements, studying industries, writing, building a skill, improving a product, or understanding a market. The person who waits for a burst of enthusiasm will be beaten by the person who returns on schedule.

A small amount of attention invested weekly into the right problem can change a life. The change may not be visible in the first week, or the tenth week, or even the first year. That is why patience is not a personality ornament. It is part of the machinery.

\section*{No Instant Feedback}

Many correct actions feel unrewarded for a long time.

Reading annual reports may feel slow. Updating a spreadsheet may feel dull. Searching unfamiliar terms may feel clumsy. Writing down one's own investment thesis may feel unnecessary when someone on the internet already sounds confident.

But these are exactly the actions that build an investor.

When entering a new field, fear often comes from unfamiliarity rather than true difficulty. The first few days are uncomfortable because every word seems new. Continue for a few more days. Read the same thing again. Search basic explanations. Write down definitions. Compare examples. Soon the strange becomes familiar, and the familiar becomes usable.

This is also how one escapes dependency. Search ability, information processing, summarizing, judgment, and writing are not gifts other people can perform on your behalf. They are muscles built through use.

A useful rule is:

\begin{quote}
When the problem becomes difficult, do not ask first. Work first.
\end{quote}

This does not mean never learning from others. It means not using others to avoid thinking. Good questions are valuable after effort. Lazy questions weaken the person asking them.

In investing, that weakness is punished.

\section*{When to Sell}

Dollar-cost averaging also simplifies the selling question, but only if the original selection was sound.

If the company is still growing and the original thesis remains intact, there may be no urgent reason to exit completely. The investor may rebalance, reduce exposure, or redirect new contributions elsewhere, but the default question should not be ``How do I escape?'' The default question should be ``Is the growth thesis still true?''

This is a more useful question because price alone is not enough. A rising price may still be justified by rising value. A falling price may be an opportunity if the business is stronger than the market believes. Or it may be a warning if the business itself is deteriorating.

The regular buyer must therefore review the thesis, not the mood.

A simple review structure can help:

\begin{enumerate}
\item Has the company's long-term growth engine changed?
\item Has the competitive position improved, weakened, or stayed the same?
\item Are the financial results confirming or contradicting the original idea?
\item Has the price become so high that future return is likely poor even if the company remains good?
\item Is there a better use for new capital?
\end{enumerate}

These questions are not a mechanical sell signal. They are a guard against sleepwalking.

\section*{The Annual Assignment}

A person who wants to become an investor should not begin by asking for stock tips. A better beginning is to build a simple research habit that lasts at least a year.

Create a table of companies. Once a month, update their prices. Add basic notes. Record what you believed before and what changed afterward. Add a dollar-cost averaging column and calculate what would have happened if you had invested a fixed amount at each interval.

Do not ask someone else how the table must look. Make one. Improve it. Break it. Rebuild it. The point is not the elegance of the spreadsheet. The point is the training of judgment.

Then add a second habit: each month, identify one more company that may have long-term growth characteristics. Search. Read. Compare. Write down why it might be worth studying and what would prove the idea wrong.

After twelve months, the investor will not merely have a table. He will have a record of attention. He will have watched prices move without needing to react to every movement. He will have experienced the gap between opinion and evidence. He will have become less frightened by unfamiliar documents, less dependent on other people's conclusions, and more capable of asking precise questions.

A year passes quickly whether the work is done or not. The difference is that one version of the year leaves behind only memory. The other leaves behind skill.

\section*{A Pocket Rule}

The simplest safe investment strategy is not ``buy anything regularly.''

It is:

\begin{quote}
After serious study, regularly invest a fixed amount into assets with genuine long-term growth, while continuing to review the growth thesis.
\end{quote}

This sentence contains the whole method.

The fixed schedule protects against emotional timing. The fixed amount protects against reckless sizing. The growth requirement protects against buying mediocrity forever. The review protects against blind persistence.

The method is simple. That is why people underestimate it. They would rather search for a secret that feels more advanced. But in wealth freedom, the advanced person is often the one who can actually execute the simple truth for years.

Buy low, sell high. Improve a little every week. Invest attention before money. Do homework before action. Choose the right thing before optimizing speed. Repeat long enough for compounding to appear.

These are public secrets. They become private advantages only when practiced.
